% ---------------------------------------------------------------------------
%  How to read this book
% ---------------------------------------------------------------------------
\chapter*{How to Read This Book}
\addcontentsline{toc}{chapter}{How to Read This Book}
\markboth{How to Read This Book}{How to Read This Book}

\section*{Two tracks, one runtime}

This book contains two courses that share a single platform. Part~II is the
\textbf{Computer Architecture track}, written entirely in \riscv assembly.
Part~III is the \textbf{Computer Programming track}, written entirely in C. The
two parts are independent: neither requires you to have read the other, and they
deliberately solve some of the same problems so that they can be compared. If you
are following only one course, read the shared foundation in Part~I, then your
own track, and consult the appendices as needed. Part~IV is a short closing
chapter that brings the two languages into contact; it rewards readers who have
seen both, but it is optional.

\begin{center}
\begin{tikzpicture}[node distance=8mm and 12mm,
  every node/.style={font=\sffamily\small},
  box/.style={rounded corners=2pt,draw=prgsteel,fill=prgpanel,minimum height=9mm,minimum width=34mm,align=center},
  acc/.style={rounded corners=2pt,draw=prggreen,fill=prggreen!8,minimum height=9mm,minimum width=34mm,align=center},
  cbox/.style={rounded corners=2pt,draw=prgamber,fill=prgamber!8,minimum height=9mm,minimum width=34mm,align=center},
  ar/.style={-{Stealth[length=2mm]},draw=prgink!60}]
  \node[box] (found) {Part~I\\Foundations};
  \node[acc,below left=of found] (asm) {Part~II\\Assembly track};
  \node[cbox,below right=of found] (c) {Part~III\\C track};
  \node[box,below=18mm of found] (join) {Part~IV\\Assembly meets C};
  \draw[ar] (found) -- (asm);
  \draw[ar] (found) -- (c);
  \draw[ar] (asm) -- (join);
  \draw[ar] (c) -- (join);
\end{tikzpicture}
\end{center}

\section*{The recurring boxes}

Throughout the book a small set of coloured boxes carries specific kinds of
information. Learning to recognise them at a glance will speed up your reading.

\begin{prgconcept}
A \textbf{Concept} box states an idea you should be able to explain in your own
words before moving on. Concepts are the things an exam, or a colleague, will ask
you about.
\end{prgconcept}

\begin{prgtry}
A \textbf{Try It Yourself} box is a small, concrete task. The book is built on the
principle that you learn by changing one thing and observing the result, so these
boxes appear often and are meant to be done, not merely read.
\end{prgtry}

\begin{prgcheck}
A \textbf{Checkpoint} box lists what should be true if the preceding section
worked. If a checkpoint fails, stop and fix it before continuing; later sections
assume the earlier ones succeeded.
\end{prgcheck}

\begin{prgwarn}
A \textbf{Watch Out} box flags a mistake that beginners reliably make. Reading
these will save you hours of confusion at the workbench.
\end{prgwarn}

\begin{prgaside}
A \textbf{Historical Note} box steps back to place an idea in its context---why a
convention exists, where a technique came from, or what an old machine did. These
are enrichment; skipping them costs you nothing technical.
\end{prgaside}

\section*{Both targets, every time}

Every program in this book is designed to run on two targets without changing a
single line of game code: Espressif's \textbf{QEMU} emulator on your laptop, and
the physical \textbf{ESP32\nobreakdash-C6} board. When a section gives build commands, it
gives them for both. The QEMU target uses an ESP32\nobreakdash-C3 emulator profile because
that is the \riscv graphics path Espressif maintains, while the hardware target
is the ESP32\nobreakdash-C6; the \riscv calling convention and the \prg interface are
identical across the two \citep{prg32qemu}. Wherever a feature exists only on real
hardware---physical buttons, the buzzer, Wi\nobreakdash-Fi cartridge upload---the text says
so plainly.

\section*{What you need}

To follow the book you need a computer running Windows, Linux, or macOS, and the
free, open-source toolchain described in Appendix~F. To follow the optional
hardware exercises you also need an ESP32\nobreakdash-C6 development board and a small
collection of inexpensive parts listed in Appendix~E. Nothing else in this book
is proprietary or paid.
